% \iffalse meta-comment
%
%% regulatory-attachments.dtx
%% Copyright 2024-2026 E. Nijenhuis
%
% This work may be distributed and/or modified under the
% conditions of the LaTeX Project Public License, either version 1.3c
% of this license or (at your option) any later version.
% The latest version of this license is in
% http://www.latex-project.org/lppl.txt
% and version 1.3c or later is part of all distributions of LaTeX
% version 2005/12/01 or later.
%
% This work has the LPPL maintenance status ‘maintained’.
%
% The Current Maintainer of this work is E. Nijenhuis.
%
% This work consists of the files listed in the meta-comment of
% regulatory-struct.dtx.
%
% \fi
%
% \iffalse
%<*driver>
\ProvidesFile{regulatory-attachments.dtx}
%</driver>
%<package>\NeedsTeXFormat{LaTeX2e}
%<package>\ProvidesPackage{regulatory-attachments}
%<package>    [2026/09/05 0.0.5 Xerdi's Regulatory Package (Attachments)]
%
%<*driver>
\documentclass[10pt,english]{ltxdoc}
%! suppress = InclusionLoop
\usepackage{regulatory}
\usepackage{tabularx}
\usepackage[english,dutch]{babel}
%% regulatory-preamble.tex
%% Copyright 2024 E. Nijenhuis
%
% This work may be distributed and/or modified under the
% conditions of the LaTeX Project Public License, either version 1.3c
% of this license or (at your option) any later version.
% The latest version of this license is in
% http://www.latex-project.org/lppl.txt
% and version 1.3c or later is part of all distributions of LaTeX
% version 2005/12/01 or later.
%
% This work has the LPPL maintenance status ‘maintained’.
%
% The Current Maintainer of this work is E. Nijenhuis.
%
% This work consists of the files regulatory.ins,
% regulatory-struct.dtx, regulatory-ref.dtx,
% regulatory-defs.dtx, regulatory-attachments.dtx,
% regulatory-md.dtx,
% regulatory-html.dtx,
% regulatory-english.def, regulatory-dutch.def,
% regulatory.tex, regulatory-preamble.tex,
% regulatory-nl.tex, regulatory-en.tex,
% example1.bib, example2.bib,
% example1-nl.tex, example2-nl.tex,
% example1-en.tex, example2-en.tex,
% md-example.tex, example.md,
% and the derived files regulatory.sty, regulatory-struct.sty,
% regulatory-defs.sty, regulatory-ref.sty,
% regulatory-attachments.sty, regulatory-md.sty
% and regulatory.4ht.
\usepackage{gitinfo-lua}
\usepackage{listings}
\usepackage{adjustbox}
\usepackage{tikz}
\usetikzlibrary{arrows.meta}

\usepackage{fontspec}
\usepackage{textcomp}

\def\projecturl{https://github.com/Xerdi/regulatory}

\definecolor{greycolor}{HTML}{BFBFBF}
\definecolor{primary}{HTML}{174A67}
\definecolor{darkcolor}{HTML}{091D28}
\definecolor{greencolor}{HTML}{176734}
\colorlet{artlabel}{primary}
\colorlet{green}{greencolor}

\lstdefinestyle{tex}{%
    language=[LaTeX]TeX,
    keywordsprefix={\\},
    alsoletter={\\},
    morekeywords={article,para,textfill,masterdocument,refdocument,loadglsdefs,printdefs,setmiddleconjunction,setlastconjunction,setrangeconjunction,setconjunction,markdownInput,glssetwidest,describe}
}

\lstdefinestyle{bash}{%
    language=bash,
    morekeywords={pdflatex,lualatex,latexmk,makeglossaries,bib2gls}
}

\lstdefinelanguage{Markdown}[LaTeX]{TeX}{%
    sensitive=false,
    morecomment=[l][\bfseries\ttfamily\color{green}]{\#},
    morecomment=[s][\bfseries\ttfamily\color{purple}]{:\ \ }{\ },
    morestring=[s][\color{black}]{\{}{\}},
    morekeywords={article,para,textfill,masterdocument,refdocument,loadglsdefs,printdefs,setmiddleconjunction,setlastconjunction,setrangeconjunction,setconjunction,markdownInput,glssetwidest,describe},
    keywordsprefix={\\},
    alsoletter={\\},
}

\lstdefinestyle{md}{%
    language=Markdown
}

\lstdefinelanguage{BibTeX}{%
    keywords={%
        article,book,collectedbook,conference,electronic,entry,ieeetranbstctl,%
        inbook,incollectedbook,incollection,injournal,inproceedings,%
        manual,mastersthesis,misc,patent,periodical,phdthesis,preamble,%
        proceedings,standard,string,techreport,unpublished%
    },
    keywordsprefix={@},
    comment=[l][\itshape]{@comment},
    sensitive=false,
}

\lstdefinestyle{bib}{%
    language=BibTeX,
    morestring=[s][\bfseries\ttfamily\color{purple}]{\ \ \ \ }{\ },
    morecomment=[n][\bfseries\ttfamily\color{green}]{\ \{}{\}},
}

\lstset{
    title=\lstname,
    basicstyle=\ttfamily,
    numberstyle=\ttfamily,
    numbersep=5pt,
    rulecolor=\color{greycolor},
    stringstyle=\color{pink},
    keywordstyle=\color{primary},
    commentstyle=\itshape\color{darkcolor},
    breaklines=true,
    captionpos=b,
    tabsize=4,
    showtabs=true
}

\newcommand\package{\texttt}
\newcommand\file{\texttt}
\newcommand\option{\texttt}

% Reads test/conformance.tex, which test/conformance.sh generates and which is
% kept in the repository: veraPDF is not available everywhere this manual is
% built, so the verdicts travel with the sources.
\newcommand\conformanceversion{?}
\newcommand\conformanceimage{?}
% A digest carries no break points of its own and does not fit a line, so this
% offers one after every character. The argument is expanded first, otherwise the
% loop sees one macro instead of the characters it stands for.
\newcommand\anywherebreakend{}
\def\anywherebreakloop#1{%
    \ifx#1\anywherebreakend\else#1\allowbreak\expandafter\anywherebreakloop\fi}
\newcommand\anywherebreak[1]{\anywherebreakloop#1\anywherebreakend}
\newcommand\conformanceimagetext{%
    \expandafter\anywherebreak\expandafter{\conformanceimage}}

% \ignorespaces, because this macro typesets nothing while its line in the
% generated file still ends in a newline, and that space would otherwise open
% the first cell of the table.
\newcommand\conformancevalidator[2]{%
    \gdef\conformanceversion{#1}\gdef\conformanceimage{#2}\ignorespaces}
\newcommand\conformancerow[6]{%
    \texttt{#1} & \texttt{#2} & \texttt{#3} & \texttt{#4} & \texttt{#5}%
    \ifthenelse{\equal{#6}{}}{}{ \footnotesize(#6)} \\}

% Points at the resulting PDF of an example, which is attached to this manual.
% Its argument is the label of the artifact the example was declared with, and it
% falls back to plain text whenever the PDF isn't there.
\newcommand\exampleresult[1]{%
    \par\noindent
    \translation{The result of this example is}{Het resultaat van dit voorbeeld is}
    \documentattachment{#1}{\file{\artifactfilename{#1}}}.\par}

\def\packagedate{\gitdate}
\def\packageversion{\gitversion}

%\def\cmd{\lstinline[style=tex]}
\def\cmditem#1{\item[\ttfamily\textbackslash{}#1]}

\newcommand\Vtextvisiblespace[1][.3em]{%
  \mbox{\kern.06em\vrule height.3ex}%
  \vbox{\hrule width#1}%
  \hbox{\vrule height.3ex}}

% Package zref-clever loads its language files lazily, which would otherwise
% happen inside a \DocInput region, where the percent sign is ignored and the
% comments of those files end up being parsed as keys. Touching every language
% of this manual schedules them at the start of the document instead, while the
% percent sign still is a comment character.
\zcsetup{lang=english}
\zcsetup{lang=dutch}
\zcsetup{lang=current}

\newcommand\translation[2]{#1}
\begin{document}
    \selectlanguage{english}
    \DocInput{regulatory-attachments.dtx}
\end{document}
%</driver>
% \fi
%
% \subsection{\texorpdfstring{\package{regulatory-attachments}}{regulatory-attachments}}
% \setcounter{CodelineNo}{0}
%
% \subsubsection{\translation{Prerequisites}{Vereisten}}
%
% A reference to another document is a reference first, so this module builds on \package{regulatory-ref},
% which loads \package{regulatory-struct} in turn. The definitions of another document land in a glossary
% of their own, which is why \package{regulatory-defs} is needed as well. Both \package{xr-hyper} and \package{zref-xr} are already
% loaded by \package{regulatory-struct}, since they have to be loaded before \package{hyperref}.
% Options are handed over to \package{regulatory-struct}, which administers them for the whole bundle.
% \iffalse
%<*package>
% \fi
%    \begin{macrocode}
\@ifpackageloaded{regulatory-struct}{}{%
    \DeclareOption*{\PassOptionsToPackage{\CurrentOption}{regulatory-struct}}%
    \ProcessOptions\relax
}
\RequirePackage{regulatory-ref}
\RequirePackage{regulatory-defs}

\DeclareTranslationFallback{of the}{of the}
\DeclareTranslationFallback{see}{see}
%    \end{macrocode}
%
% Definitions of other documents are kept in a glossary type of their own, so they don't get mixed up with
% the definitions of the document itself.
%    \begin{macrocode}
\newglossary*{externals}{Externals}
%    \end{macrocode}
%
% \subsubsection{\translation{Artifacts}{Artefacten}}
%
% \begin{macro}{\newartifact}
% An artifact holds everything known about another document: where its file is, how it is referred to, and
% what a PDF viewer should say about it once it is attached. Every artifact is a key family of its own,
% named after \meta{name}.
%    \begin{macrocode}
\newcommand\artifact@footnote@noop[1]{#1}
\newcommand\artifact@def@label@default[1]{~(\GetTranslation{see} #1)}

\newcommand\newartifact[2]{%
    \pgfkeys{
        /artifact/#1/.is family,
        /artifact/#1,
        name/.initial = #1,
        filename/.initial = #1.pdf,
        path/.initial = ./,
        ref label/.initial=,
        def label/.initial=\artifact@def@label@default,
        footnote/.initial = true,
        footnote label/.initial=\artifact@footnote@noop,
        url/.initial=,
        referred/.initial = false,
        defs/.initial=,
        author/.initial=<author>,
        subject/.initial=<subject>,
        description/.initial=<description>
    }
    \pgfkeys{/artifact/#1,#2}
}
%    \end{macrocode}
% \end{macro}
%
% The accessors of an artifact, of which the full name and the full filename are the ones prefixed with the
% path of the artifact.
%    \begin{macrocode}
\newcommand\@aget[2]{\pgfkeysvalueof{/artifact/#1/#2}}
\newcommand\artifact@setfootnotelabel[1]{\pgfkeysgetvalue{/artifact/#1/footnote label}{\artifactfootnotelabel}}
\newcommand\artifactname[1]{\@aget{#1}{name}}
\newcommand\artifactfullname[1]{\@aget{#1}{path}\@aget{#1}{name}}
\newcommand\artifactfilename[1]{\@aget{#1}{filename}}
\newcommand\artifactfullfilename[1]{\@aget{#1}{path}\@aget{#1}{filename}}
\newcommand\artifactfootnote[1]{\@aget{#1}{footnote}}
\newcommand\artifacturl[1]{\@aget{#1}{url}}
\newcommand\artifactreferred[1]{\@aget{#1}{referred}}
\newcommand\artifactdefs[1]{\@aget{#1}{defs}}
\newcommand\artifactauthor[1]{\@aget{#1}{author}}
\newcommand\artifactsubject[1]{\@aget{#1}{subject}}
\newcommand\artifactdescription[1]{\@aget{#1}{description}}
%    \end{macrocode}
%
% \subsubsection{\translation{Embedding files}{Bestanden invoegen}}
%
% \begin{macro}{\regulatory@embed}
% Embeds the file of \#1 under the name \#2, as an object named \#3 with description \#4. The file is
% registered as an associated file of the document with relationship \texttt{/Source}, and it is added to
% the embedded files of the catalog, so that PDF viewers list it as an attachment.
%    \begin{macrocode}
\ExplSyntaxOn
\str_new:N \l__regulatory_desc_str
\str_new:N \l__regulatory_ext_str
%    \end{macrocode}
%
% PDF/A-3 requires the MIME type of every embedded file, and prescribes
% \texttt{application/octet-stream} whenever it is unknown. \package{l3pdf} looks the type up by
% extension and only warns when it finds none, so the prescribed fallback is filled in here.
%    \begin{macrocode}
\cs_new_protected:Npn \__regulatory_mimetype:n #1
  {
    \file_parse_full_name:nNNN {#1}
      \l_tmpa_str \l_tmpb_str \l__regulatory_ext_str
    \prop_if_in:NVF \g_pdffile_mimetypes_prop \l__regulatory_ext_str
      {
        \pdfdict_put:nne {l_pdffile}{Subtype}
          { \pdf_name_from_unicode_e:n {application/octet-stream} }
      }
  }

\cs_new_protected:Npn \__regulatory_embed:nnnn #1#2#3#4
  {
    \group_begin:
      \pdfdict_put:nnn {l_pdffile/Filespec}{AFRelationship}{/Source}
      \__regulatory_mimetype:n {#1}
      \tl_if_empty:nF {#4}
        {
          \pdf_string_from_unicode:nnN {utf16/string}{#4}\l__regulatory_desc_str
          \pdfdict_put:nne {l_pdffile/Filespec}{Desc}{\l__regulatory_desc_str}
        }
      \pdffile_embed_file:nnn {#1}{#2}{regulatory/attach/#3}
      \pdfmanagement_add:nne {Catalog/Names/EmbeddedFiles}
        {#3}{\pdf_object_ref:n{regulatory/attach/#3}}
      \pdfmanagement_add:nne {Catalog}{AF}
        {\pdf_object_ref:n{regulatory/attach/#3}}
    \group_end:
  }


\cs_generate_variant:Nn \__regulatory_embed:nnnn {eeen}
\cs_new_protected:Npn \regulatory@embed #1#2#3#4
  { \__regulatory_embed:eeen {#1}{#2}{#3}{#4} }
%    \end{macrocode}
% \end{macro}
%
% \begin{macro}{\regulatory@attachmentlink}
% Turns \#2 into a link that opens the embedded file named \#1 in a new window.
%    \begin{macrocode}
\NewDocumentCommand \regulatory@attachmentlink { m m }
  {
    \pdfannot_link_begin:nnw {user}
      { /Subtype/Link/F~4/A<</S/GoToE/NewWindow~true/T<</R/C/N~(#1)>>>> }
    #2
    \pdfannot_link_end:n {user}
  }
%    \end{macrocode}
% \end{macro}
%
% \begin{macro}{\regulatory@ifembed}
% Embedding a file is only possible while the PDF management of \LaTeX{} is active, which a document
% activates with \cmd{\DocumentMetadata} in front of its \cmd{\documentclass}. Without it,
% the file embedding commands of \package{l3pdf} do not even exist, so every attachment falls back to its plain
% text and the document is told once what it is missing.
%    \begin{macrocode}
\cs_if_exist:NTF \pdffile_embed_file:nnn
  { \cs_new_eq:NN \regulatory@ifembed \use_i:nn }
  { \cs_new_eq:NN \regulatory@ifembed \use_ii:nn }
\ExplSyntaxOff

\regulatory@ifembed{}{%
    \PackageWarning{regulatory}{%
        Attachments need the PDF management of LaTeX.\MessageBreak
        Put \string\DocumentMetadata{} before \string\documentclass\MessageBreak
        to activate it. Attachments are typeset as plain text%
    }%
}
%    \end{macrocode}
% \end{macro}
%
% \begin{macro}{\regulatory@embedonce}
% A document referred to more than once is embedded only the first time, of which the object name is kept
% for all following links.
%    \begin{macrocode}
\newcommand\regulatory@embedonce[1]{%
    \expandafter\ifx\csname regulatory@embedded@#1\endcsname\relax
        \regulatory@embed
            {\artifactfullfilename{#1}}%
            {\artifactfilename{#1}}%
            {regulatory-#1}%
            {\artifactdescription{#1}}%
        \expandafter\xdef\csname regulatory@embedded@#1\endcsname{regulatory-#1}%
    \fi
}
%    \end{macrocode}
% \end{macro}
%
% \begin{macro}{\documentattachment}
% Attaches the file of artifact \meta{label} and prints \meta{link text} as the link to it. A missing file
% is no reason to fail; its link text is printed verbatim instead.
%    \begin{macrocode}
\newcommand\documentattachment[2]{%
    \def\@tmpfile{\artifactfullfilename{#1}}%
    \regulatory@ifembed{%
        \IfFileExists{\@tmpfile}{%
            \regulatory@embedonce{#1}%
            \expandafter\expandafter\expandafter\regulatory@attachmentlink
                \expandafter\expandafter\expandafter
                {\csname regulatory@embedded@#1\endcsname}{#2}%
        }%
        {\texttt{#2}}%
    }{\texttt{#2}}%
}
%    \end{macrocode}
% \end{macro}
%
% \begin{macro}{\attachdocument}
% Attaches any file by its \meta{filename}, without an artifact of its own, with \meta{description} as the
% description of the attachment.
%    \begin{macrocode}
\newcommand\attachdocument[3][]{%
    \regulatory@ifembed{%
        \IfFileExists{#2}{%
            \expandafter\ifx\csname regulatory@plain@#2\endcsname\relax
                \regulatory@embed {#2}{#2}{regulatory-#2}{#1}%
                \expandafter\xdef\csname regulatory@plain@#2\endcsname{regulatory-#2}%
            \fi
            \expandafter\expandafter\expandafter\regulatory@attachmentlink
                \expandafter\expandafter\expandafter
                {\csname regulatory@plain@#2\endcsname}{#3}%
        }%
        {\texttt{#3}}%
    }{\texttt{#3}}%
}
%    \end{macrocode}
% \end{macro}
%
% \subsubsection{\translation{Mentioning the source}{Bronvermelding}}
%
% \begin{macro}{\documentlabel}
% How a reference mentions the document it belongs to, which defaults to the subject of the artifact.
%    \begin{macrocode}
\newcommand\documentlabel[1]{%
    \def\@@label{\@aget{#1}{ref label}}%
    \ifthenelse{\equal{\@@label}{}}{%
        \GetTranslation{of the} \artifactsubject{#1}%
    }{\@@label}%
}
%    \end{macrocode}
% \end{macro}
%
% \begin{macro}{\documentfootnote}
% Places the footnote of the first occurrence of a reference or a definition of another document, with the
% document itself attached to it. The artifact is marked as referred, so that following occurrences don't
% repeat the footnote.
%    \begin{macrocode}
\newcommand\documentfootnote[2][]{%
    \pgfkeysifdefined{/artifact/#2/referred}{%
        \pgfkeys{/artifact/#2,referred=true}%
        \artifact@setfootnotelabel{#2}%
        \ifthenelse{\equal{#1}{}}{%
            \def\artifact@description{\artifactsubject{#2}}%
        }{%
            \def\artifact@description{#1}%
        }%
        \footnote{\artifactfootnotelabel{\documentattachment{#2}{\artifact@description}}}%
    }{\PackageWarning{regulatory}{Artifact does not exists '#2'}\footnote{DOCUMENT NOT FOUND '#2'}}
}
%    \end{macrocode}
% \end{macro}
%
% \begin{macro}{\aref@postlink@setup}
% This is where a complete reference of \package{regulatory-ref} gets its source mentioned. Labels of the
% document itself carry no \texttt{externaldocument} property, in which case nothing is added at all.
%    \begin{macrocode}
\renewcommand\aref@postlink@setup[1]{%
    \zref@ifrefcontainsprop{#1}{externaldocument}{%
        \filename@parse{\expanded{\zref@extract{#1}{externaldocument}}}%
        \edef\@aref@filelabel{\expanded{\filename@base}}%
        \def\@aref@postlink{%
            , \documentlabel{\@aref@filelabel}%
            \ifthenelse{\equal{\artifactfootnote{\@aref@filelabel}}{true}\and\equal{\artifactreferred{\@aref@filelabel}}{false}}{%
                \documentfootnote{\@aref@filelabel}%
            }{}%
        }%
    }{%
        \def\@aref@postlink{}%
    }%
}
%    \end{macrocode}
% \end{macro}
%
% \subsubsection{\translation{Linking documents}{Documenten koppelen}}
%
% \begin{macro}{\loadextdefs}
% Loads the definitions of \meta{src} under the \texttt{externals} type and the category \meta{category},
% which is the name of the artifact they belong to. Every first occurrence of such a definition mentions its
% source, unless the document was already referred to.
%    \begin{macrocode}
\newcommand\loadextdefs[2][externals]{%
    \ifregulatory@bibtogls%
        \GlsXtrLoadResources[type=externals,src={#2},sort={nl-NL},category={#1}]%
    \else%
        \loadglsentries[externals]{#2}%
    \fi%
    \glsdefpostlink{#1}{%
        \ifthenelse{\equal{\artifactreferred{#1}}{true}}{}{%
            \@aget{#1}{def label}{%
                \artifactsubject{#1}%
                \documentfootnote{#1}%
            }%
        }%
    }%
}
%    \end{macrocode}
% \end{macro}
%
% \begin{macro}{\refdocument}
% Declares another \package{regulatory} document to refer to, of which the labels are read with
% \cmd{\zexternaldocument} of \package{zref-xr} and prefixed with \meta{prefix}.
%    \begin{macrocode}
\newcommand\refdocument[3][ext-]{%
    \newartifact{#2}{footnote=false,#3}%
    \zexternaldocument[#1]{\artifactfullname{#2}}[#2]%
}
%    \end{macrocode}
% \end{macro}
%
% \begin{macro}{\masterdocument}
% Like \cmd{\refdocument}, but with the definitions of the other document loaded as well and with the
% document itself attached at the first occurrence of a reference or a definition. With \package{bib2gls},
% the hyperlinks of those definitions are aimed at the other document.
%    \begin{macrocode}
\newcommand\masterdocument[3][ext-]{%
    \newartifact{#2}{#3}%
    \zexternaldocument[#1]{\artifactfullname{#2}}[#2]%
    \ifthenelse{\equal{\artifactdefs{#2}}{}}{%
        \PackageWarning{regulatory}{No definitions given to the master document. \the\gls commands will therefore not work.}%
    }{%
        \loadextdefs[#2]{\artifactdefs{#2}}%
        \ifregulatory@bibtogls%
            \glssetcategoryattribute{#2}{targeturl}{\artifactfilename{#2}}%
            \glssetcategoryattribute{#2}{targetname}{\glolinkprefix\glslabel}%
        \fi%
    }%
}
%    \end{macrocode}
% \iffalse
%</package>
% \fi
% \end{macro}
%
% \Finale
%
